
\documentclass{article}
\usepackage{colortbl}
\usepackage{makecell}
\usepackage{multirow}
\usepackage{supertabular}

\begin{document}

\newcounter{utterance}

\centering \large Interaction Transcript for game `codenames', experiment `abstract', episode 2 with reflexion{-}llama{-}3.1{-}8b{-}instruct.
\vspace{24pt}

{ \footnotesize  \setcounter{utterance}{1}
\setlength{\tabcolsep}{0pt}
\begin{supertabular}{c@{$\;$}|p{.15\linewidth}@{}p{.15\linewidth}p{.15\linewidth}p{.15\linewidth}p{.15\linewidth}p{.15\linewidth}}
    \# & $\;$A & \multicolumn{4}{c}{Game Master} & $\;\:$B\\
    \hline

    \theutterance \stepcounter{utterance}  
    & & & \multicolumn{2}{p{0.3\linewidth}}{
        \cellcolor[rgb]{0.9,0.9,0.9}{
            \makecell[{{p{\linewidth}}}]{
                \texttt{\tiny{[GM$|$GM]}}
                {'hidden': {'team': ['desirability', 'renegotiation', 'coincidence'], 'innocent': ['graciousness', 'alienation', 'inconsequence', 'seriousness', 'unification', 'optimism'], 'opponent': [], 'assassin': ['furtherance']}, 'revealed': {'team': {'team': [], 'innocent': [], 'opponent': [], 'assassin': []}, 'opponent': {'team': [], 'innocent': [], 'opponent': [], 'assassin': []}}}
            }
        }
    }
    & & \\ \\

    \theutterance \stepcounter{utterance}  
    & & \multicolumn{4}{p{0.6\linewidth}}{
        \cellcolor[rgb]{0.9,0.9,0.9}{
            \makecell[{{p{\linewidth}}}]{
                \texttt{\tiny{[P1$\langle$GM]}}
                \texttt{You have attempted this same task before and did not succeed. These are the plans you wrote after those attempts.} \\
\texttt{Trial \#1: New plan: } \\
\\ 
\texttt{In my previous attempts, I provided clues that were too specific and only related to a few words. This led to the opponent guessing some of the words, which ended my teammate's turn. I also failed to consider the remaining words for the opponent and the distractor words. Next time, I will provide a clue that is more general and relates to multiple words, while also taking into account the remaining words for the opponent and the distractor words. I will also make sure to follow the format and provide a well{-}formed and legal response. For example, I could provide a clue that relates to the theme of the words, such as "Morality" for words like shamelessness, penance, and motive.} \\
\\ 
\texttt{Trial \#2: New plan: } \\
\\ 
\texttt{In my previous attempts, I provided clues that were too specific and only related to a few words. This led to the opponent guessing some of the words, which ended my teammate's turn. I also failed to consider the remaining words for the opponent and the distractor words. Next time, I will provide a clue that is more general and relates to multiple words, while also taking into account the remaining words for the opponent and the distractor words. I will also make sure to follow the format and provide a well{-}formed and legal response.} \\
\\ 
\texttt{To achieve this, I will start by analyzing the theme and common characteristics of the remaining team words. In this case, the words are desirability, renegotiation, and coincidence. I will look for a clue that relates to a broader concept that encompasses these words, such as "Relationships" or "Interactions". I will then make sure to avoid any words that are too specific or related to a few words, and instead opt for a more general clue that can be applied to multiple words.} \\
\\ 
\texttt{Additionally, I will take into account the remaining words for the opponent and the distractor words, and make sure that my clue does not inadvertently provide information that could help the opponent guess their words. I will also be mindful of the kill words and avoid providing a clue that could lead to the game ending prematurely.} \\
\\ 
\texttt{With this new approach, I will provide a clue that is more general, yet still relevant to the remaining team words, and avoid any potential pitfalls that could lead to a loss.} \\
\\ 
\\ 
\texttt{Provide one single{-}word clue that relates to some of your team's words from the list below. You can choose to how many and to which words your clue relates to.} \\
\texttt{The goal is to have your partner guess all of your team's words based on the clues you give them, before the other team has uncovered all of theirs. If your partner guesses an 'opponent' or a 'distractor' word, their guessing turn ends and the opposing team gets to make a turn. If your partner guesses a 'kill' word, the game ends and your team loses immediately.} \\
\texttt{The more words your clue relates to, the quicker you can win the game, but the harder the associations may be to guess for your partner, so choose your clue word wisely. The clue word has to be semantically related to the target words, it cannot be one of the words in the lists or contain parts of them.} \\
\\ 
\texttt{Always give your single{-}word clue and your comma{-}separated list of related target words in the following format and make your answers as short as possible, never include any other text than is required in this form:} \\
\\ 
\texttt{CLUE: <WORD>} \\
\texttt{TARGETS: <TARGETS>} \\
\\ 
\texttt{Your team words are: desirability, renegotiation, coincidence.} \\
\texttt{Your opponent's team words are: .} \\
\texttt{Distractor words are: graciousness, alienation, inconsequence, seriousness, unification, optimism.} \\
\texttt{The kill words are: furtherance.} \\
            }
        }
    }
    & \\ \\

    \theutterance \stepcounter{utterance}  
    & \multicolumn{4}{p{0.6\linewidth}}{
        \cellcolor[rgb]{0.9,0.9,0.9}{
            \makecell[{{p{\linewidth}}}]{
                \texttt{\tiny{[P1$\rangle$GM]}}
                \texttt{CLUE: Chance} \\
\texttt{TARGETS: desirability, renegotiation, coincidence} \\
            }
        }
    }
    & & \\ \\

    \theutterance \stepcounter{utterance}  
    & & & \multicolumn{2}{p{0.3\linewidth}}{
        \cellcolor[rgb]{0.9,0.9,0.9}{
            \makecell[{{p{\linewidth}}}]{
                \texttt{\tiny{[GM$|$GM]}}
                \texttt{chance} \\
            }
        }
    }
    & & \\ \\

    \theutterance \stepcounter{utterance}  
    & & & \multicolumn{2}{p{0.3\linewidth}}{
        \cellcolor[rgb]{0.9,0.9,0.9}{
            \makecell[{{p{\linewidth}}}]{
                \texttt{\tiny{[GM$|$GM]}}
                ['desirability', 'renegotiation', 'coincidence']
            }
        }
    }
    & & \\ \\

    \theutterance \stepcounter{utterance}  
    & & & \multicolumn{2}{p{0.3\linewidth}}{
        \cellcolor[rgb]{0.9,0.9,0.9}{
            \makecell[{{p{\linewidth}}}]{
                \texttt{\tiny{[GM$|$GM]}}
                {'word': 'desirability', 'assignment': 'team'}
            }
        }
    }
    & & \\ \\

    \theutterance \stepcounter{utterance}  
    & & & \multicolumn{2}{p{0.3\linewidth}}{
        \cellcolor[rgb]{0.9,0.9,0.9}{
            \makecell[{{p{\linewidth}}}]{
                \texttt{\tiny{[GM$|$GM]}}
                {'word': 'renegotiation', 'assignment': 'team'}
            }
        }
    }
    & & \\ \\

    \theutterance \stepcounter{utterance}  
    & & & \multicolumn{2}{p{0.3\linewidth}}{
        \cellcolor[rgb]{0.9,0.9,0.9}{
            \makecell[{{p{\linewidth}}}]{
                \texttt{\tiny{[GM$|$GM]}}
                {'word': 'coincidence', 'assignment': 'team'}
            }
        }
    }
    & & \\ \\

    \theutterance \stepcounter{utterance}  
    & & & \multicolumn{2}{p{0.3\linewidth}}{
        \cellcolor[rgb]{0.9,0.9,0.9}{
            \makecell[{{p{\linewidth}}}]{
                \texttt{\tiny{[GM$|$GM]}}
                \texttt{opponent has won} \\
            }
        }
    }
    & & \\ \\

    \theutterance \stepcounter{utterance}  
    & \multicolumn{4}{p{0.6\linewidth}}{
        \cellcolor[rgb]{0.95,0.95,0.95}{
            \makecell[{{p{\linewidth}}}]{
                \texttt{\tiny{[P2$\rightarrow$P1]}}
                \texttt{Post{-}episode Reflexion (not sent to GM):} \\
\texttt{New plan: } \\
\\ 
\texttt{In my previous attempts, I provided clues that were too specific and only related to a few words. This led to the opponent guessing some of the words, which ended my teammate's turn. I also failed to consider the remaining words for the opponent and the distractor words. Next time, I will provide a clue that is more general and relates to multiple words, while also taking into account the remaining words for the opponent and the distractor words. I will also make sure to follow the format and provide a well{-}formed and legal response.} \\
\\ 
\texttt{To achieve this, I will start by analyzing the theme and common characteristics of the remaining team words. In this case, the words are desirability, renegotiation, and coincidence. I will look for a clue that relates to a broader concept that encompasses these words, such as "Relationships" or "Interactions". I will then make sure to avoid any words that are too specific or related to a few words, and instead opt for a more general clue that can be applied to multiple words.} \\
\\ 
\texttt{Additionally, I will take into account the remaining words for the opponent and the distractor words, and make sure that my clue does not inadvertently provide information that could help the opponent guess their words. I will also be mindful of the kill words and avoid providing a clue that could lead to the game ending prematurely.} \\
\\ 
\texttt{I will also consider the semantic relationships between the words, such as synonyms, antonyms, hyponyms, and hypernyms, to provide a clue that is more accurate and relevant to the remaining team words. For example, if the words are related to each other through a common concept, I will provide a clue that is related to that concept, rather than a specific word.} \\
\\ 
\texttt{With this new approach, I will provide a clue that is more general, yet still relevant to the remaining team words, and avoid any potential pitfalls that could lead to a loss.} \\
            }
        }
    }
    & & \\ \\

\end{supertabular}
}

\end{document}
